\section{Introduction}

On parallel FLORES-200 text, Qwen3's tokenizer needs $4.42$ times as many tokens for a Nepali sentence as for its English translation, and the \texttt{cl100k} tokenizer used by Phi-4, OLMo-2 and Granite-4.1 needs $4.76$ times as many (\S\ref{sec:audit}). We call this ratio the \emph{premium}. Nepali speakers pay it on every request, in price, latency and usable context.

Much of the premium has a mechanical cause. GPT-2-style pre-tokenizers split text with a regular expression whose word class \pL{} matches one or more Unicode letters~\citep{gpt2encoder}; we call it the \emph{letters-only regex}. In Unicode, Devanagari vowel signs, the virama, the anusvara and the candrabindu are combining marks (\pM{}), so the regex starts a new pre-token at each of them and byte-pair encoding (BPE) can never merge across that boundary. \citet{regmi2026vowel} identified this cause and the resulting bound on token counts. They showed that adding marks to the word class cuts Nepali from 4.78 to 1.58 tokens per word, and that 63\% of the 1{,}345 text-generation repositories they examined still use the letters-only class.

Regmi et al.\ studied tokenizers trained from scratch. Most practitioners instead adapt a model that already exists, usually by vocabulary extension: they add target-language tokens, initialise their embeddings and continue pretraining~\citep{tejaswi-etal-2024-exploring,yamaguchi2026lowres,purason-etal-2026-teaching,smith2026inplace}. What is new in this paper is that deployed-model setting. We show that the bound becomes a floor that released merge-based extensions already hit; we give a repair of a deployed model's regex that provably keeps English token ids; and we test it in controlled continued pretraining against the standard continued-BPE recipe of \citet{purason-etal-2026-teaching}, with the tokenizers released.

\paragraph{Findings.} For Llama-3 and Qwen3 the \emph{pre-token floor}, the lowest premium any merge-based extension can reach, is about $2.2$ (\S\ref{sec:floor}). Ten released extensions for six scripts end within 4\% of their floors, and our own extension under the letters-only regex stalls at it. Registering \RII{}'s learned strings as added tokens does get below the floor, but a fifth of the matches then start in the middle of a word. Our repair adds marks to the word class with two string substitutions and, together with new merges, brings Nepali to parity with English while leaving every token id of mark-free, non-Devanagari text unchanged (\S\ref{sec:method}). The model-quality result is less favourable. Under our pre-registered analysis, continued pretraining with no new tokens gave the best held-out Nepali bits per byte on both models, so our main hypothesis failed. Between the two extension recipes, the repair was non-inferior to standard extension for Llama and slightly worse for Qwen, and it was clearly worse for Llama when scored after the separator used in training (\S\ref{sec:cpt}). It needs about half as many decoding steps as standard extension for the same Nepali text.

\paragraph{Contributions.} We make four contributions:
\begin{enumerate}
  \setlength\itemsep{0pt}
  \item Evidence that the pre-token floor binds released merge-based extensions, and a test of the added-token route around it (\S\ref{sec:floor}); an audit of 24 distinct deployed tokenizers with per-tokenizer floors (\S\ref{sec:audit}).
  \item An English-preserving repair of a deployed regex with a proof, a Devanagari-only variant with a stronger guarantee, and a practical pitfall in loading repaired tokenizers (\S\ref{sec:method}).
  \item A pre-registered comparison of \RO{} (no extension), \RI{} (standard extension) and \RII{} (repair + extension) in continued pretraining of two deployed models, reported against the registered hypotheses (\S\ref{sec:cpt}).
  \item The tokenizers, code and analysis plan.\footnote{\artifactnote}
\end{enumerate}
